% Results for the BISON manuscript.
% Author: fangtoast
% Email: fangtoast@foxmail.com
% Copyright (c) 2026 fangtoast. All rights reserved.

\section{Results}

We present results in the same order as the evidence chain: forward surrogate
quality (RQ1), query-space robustness (RQ2), and solver-verified inverse design
(RQ3). Each subsection closes with what the measurement establishes and what
it leaves unresolved.

\subsection{BISON Improves Full-Band Prediction}

Table~\ref{tab:forward} shows the formal three-seed comparison. BISON obtains
\(1.0299\%\pm0.0545\%\) test MRE, compared
with \(1.3309\%\pm0.0122\%\) for the strongest non-operator baseline,
ResMLP-wide. This is a 22.6\% relative reduction. The lifted CNN and plain MLP
are less accurate at 2.2234\% and 2.4764\%, respectively. The improvement is
not explained by parameter count alone: BISON uses 389,718 parameters, less
than half the 878,646 parameters of either MLP model.

\begin{table}[t]
  \caption{Formal forward comparison on the 625-structure test split. MRE and
  training time are mean\(\pm\)standard deviation over three seeds; latency and
  peak memory are means across seeds.}
  \label{tab:forward}
  \centering
  \small
  \resizebox{\columnwidth}{!}{%
  \begin{tabular}{lrrrr}
    \toprule
    Model & MRE (\%) $\downarrow$ & Params & Train (s) & ms/sample \\
    \midrule
    \method{} & $\mathbf{1.0299\pm0.0545}$ & 389,718 & 272.71 & 0.5267 \\
    ResMLP-wide & $1.3309\pm0.0122$ & 878,646 & 80.96 & 0.0064 \\
    Lifted CNN & $2.2234\pm0.0366$ & 5,458 & 39.97 & 0.0102 \\
    MLP & $2.4764\pm0.0241$ & 878,646 & 65.66 & 0.0074 \\
    \bottomrule
  \end{tabular}}
\end{table}

BISON's best per-band result is band 4 (0.9028\% MRE), while band 10 is the
hardest (1.2765\%). Thus, the gain is distributed across the band matrix rather
than arising from a single easy output. The cost is higher training and
inference time than the small dense baselines, although sub-millisecond
inference remains negligible relative to eigensolver replay.

RQ1 is therefore answered within the fixed dataset and metric: the implemented
spectral operator gives the lowest three-seed held-out full-band MRE among the
four tested models. This comparison does not isolate which BandBlock component
causes the gain, nor does it establish resolution or lattice transfer.

\subsection{Query Interpolation Works Better than Extrapolation}

The FNO query model records \(1.1004\%\pm0.0218\%\) MRE on unseen structures
over all sampled wave vectors and \(1.0470\%\pm0.0240\%\) on withheld
interpolation queries. On a held-out path segment, its error rises to
\(3.6747\%\pm0.1730\%\) and its band-order violation rate rises from 2.16\%
under interpolation to 9.22\%. This is an important negative result: physical
coordinates allow flexible queries, but they do not by themselves guarantee
reliable extrapolation.

\begin{table}[t]
  \caption{Three-seed query prediction. Values are MRE (\%) and band-order
  violation rate (\%).}
  \label{tab:query}
  \centering
  \small
  \resizebox{\columnwidth}{!}{%
  \begin{tabular}{lrrrrrr}
    \toprule
     & \multicolumn{2}{c}{All-$k$} & \multicolumn{2}{c}{Interpolation}
     & \multicolumn{2}{c}{Held segment} \\
    Model & MRE & Viol. & MRE & Viol. & MRE & Viol. \\
    \midrule
    FNO query & \textbf{1.1004} & 2.96 & \textbf{1.0470} & 2.16 & \textbf{3.6747} & 9.22 \\
    DeepONet & 1.7937 & 1.35 & 1.6308 & \textbf{0.41} & 3.7604 & \textbf{0.99} \\
    MLP query & 2.1800 & \textbf{0.83} & 2.2079 & 0.42 & 4.2548 & 1.00 \\
    \bottomrule
  \end{tabular}}
\end{table}

DeepONet has lower ordering-violation rates despite higher MRE, showing that
point accuracy and physical ordering are distinct. For this reason the query
network is retained as an audit signal, not treated as an unquestionable
second oracle.

RQ2 has a mixed answer. The coordinate-conditioned FNO interpolates held-out
points accurately, but neither the coordinate input nor the operator backbone
prevents a large error increase on an unseen path segment. The negative result
supports using the query model as a distribution-shift diagnostic rather than
as a substitute for MPB.

\subsection{No-Expert Inverse Design Survives MPB Replay}

On all 625 test targets, the one-mode generator reaches 2.1434\% BISON MRE and
2.2614\% real-MPB MRE (Table~\ref{tab:inverse}). The resulting
surrogate--replay gap is 0.118 percentage points. Its replay distribution has
median 2.2196\%, 10th percentile 1.7264\%, and 90th percentile 2.8262\%; the
95\% bootstrap interval for its mean is [2.2263, 2.2965]\%.

\begin{table}[t]
  \caption{Seed-0 inverse evaluation on the complete 625-target test split.
  Both methods use $K=128$ and Top-5 MPB replay. CI is the 95\% bootstrap
  interval of mean MPB MRE. Gap is MPB minus BISON MRE in percentage points.}
  \label{tab:inverse}
  \centering
  \small
  \resizebox{\columnwidth}{!}{%
  \begin{tabular}{lrrrrr}
    \toprule
    Generator & BISON & MPB (\%) & 95\% CI & Gap (pp) & Eff. $K$ \\
    \midrule
    Single mode (ours) & \textbf{2.1434} & \textbf{2.2614} & [2.2263,2.2965] & \textbf{0.1180} & \textbf{112.02} \\
    Four modes & 2.2477 & 2.3916 & [2.3485,2.4341] & 0.1440 & 109.24 \\
    \bottomrule
  \end{tabular}}
\end{table}

Under the identical candidate and solver budget, the four-mode baseline has
2.3916\% MPB MRE. The paired difference (four-mode minus one-mode) is 0.1302
percentage points, with bootstrap interval [0.0905, 0.1705] and two-sided sign
test \(p=0.0010\). The result supports the narrower conclusion that expert-like
mode routing was unnecessary in this seed-0 protocol. It does not establish a
universal architectural advantage or cross-seed stability.

RQ3 is answered only at the strength of this protocol: for the complete
625-target split and one training seed, BISON-ranked single-mode candidates
retain a statistically consistent paired advantage after the same five MPB
replays per target. The solver gap is nonzero, so the result validates the
budgeted workflow rather than certifying BISON as an exact physical model.
Figure~\ref{fig:inverse-examples} complements the aggregate result with four
paired examples selected jointly from the same complete replay.

\begin{figure*}[t]
  \centering
  % Canonical editable source: docs/.../inverse/paired_main4.svg
  \includegraphics[width=\textwidth]{../../docs/reports/assets/phc2d_project_full_introduction_20260724/inverse/paired_main4.pdf}
  \caption{Paired qualitative comparison between the single-mode no-expert
  generator and the four-mode mixture baseline in the seed-0, 625-target
  replay. Rows show the jointly selected best target, median target, largest
  between-method disagreement, and worst target. For each method, the figure
  displays the generated structure and the target, BISON-surrogate, and
  real-MPB band curves; the header reports MPB MRE, surrogate MRE, and their
  signed gap. These cases expose both ordinary paired behavior and difficult
  targets, but remain qualitative single-seed evidence.}
  \Description{Four rows compare single-mode and four-mode inverse-design
  candidates for jointly selected best, median, largest-disagreement, and
  worst targets. Each method column includes a generated binary unit cell,
  target and predicted band curves, MPB replay curves, and error summaries.}
  \label{fig:inverse-examples}
\end{figure*}

\subsection{Runtime and Evidence Boundary}

For the one-mode method, generating candidates for all 625 targets takes
1.408~s in total (2.25~ms/target), while BISON ranking takes 221.333~s
(354.13~ms/target). The effective pool contains 112.02 candidates on average
after binary deduplication. The query audit returns 2.1254\% MRE and differs
from BISON by 0.8457 percentage points under its own diagnostic definition;
this disagreement is retained rather than hidden.

The forward table is a formal three-seed result. The inverse table is a
complete-test but single-seed result. Earlier retrieval and direct-inverse
numbers, field/PINO residuals, and historical seed-2 comparisons are not used
as current inverse claims because they lack the same accepted replay protocol.
